\begin{helper}
Let \(W\) be a finite type, let \(H\) be a simple graph on \(W\), and let
\(s,k\) be natural numbers.  If \(H\) has no \(s\)-clique and \(H^{c}\) has no
\(k\)-clique, then
\[
  |W| < \mathrm{RamseyNumber}(s,k).
\]
\end{helper}
\begin{proof}
By \(\noderef{RamseyPropertyCounterexampleTransport}\), the two hypotheses on
\(H\) imply that
\(\mathrm{RamseyProperty}(s,k,|W|)\) is false.  Let
\[
  A=\{n\in\mathbb{N}:\mathrm{RamseyProperty}(s,k,n)\}.
\]
By \(\noderef{RamseyPropertyNonempty}\), the set \(A\) is nonempty.  By
\(\noderef{RamseyNumber}\), \(\mathrm{RamseyNumber}(s,k)=\inf A\), where the
infimum is the natural-number least element operation.

Suppose for contradiction that \(|W|<\mathrm{RamseyNumber}(s,k)\) is false.
Then \(\inf A\le |W|\).  Since \(A\) is nonempty, the natural-number least
element \(\inf A\) belongs to \(A\), so
\(\mathrm{RamseyProperty}(s,k,\inf A)\) holds.  Applying the monotonicity
theorem \(\noderef{RamseyPropertyMonotone}\) to the inequality
\(\inf A\le |W|\), we obtain \(\mathrm{RamseyProperty}(s,k,|W|)\).  This
contradicts the counterexample statement from
\(\noderef{RamseyPropertyCounterexampleTransport}\).

Therefore the supposition was impossible, and \(|W|<\mathrm{RamseyNumber}(s,k)\).
\end{proof}
